% Generated draft: requires mathematical and source review.
\documentclass[11pt]{article}
\usepackage[margin=25mm]{geometry}
\usepackage{fontspec}
\setmainfont{Latin Modern Roman}
\usepackage{amsmath,amssymb,mathtools,unicode-math}
\setmathfont{Latin Modern Math}
\usepackage{enumitem,needspace,xurl,hyperref}
\setlength{\parindent}{0pt}
\setlength{\parskip}{5pt}
\setlength{\emergencystretch}{3em}
\newcommand{\N}{\mathbb{N}}
\newcommand{\Z}{\mathbb{Z}}
\newcommand{\Q}{\mathbb{Q}}
\newcommand{\R}{\mathbb{R}}
\newcommand{\C}{\mathbb{C}}
\newcommand{\F}{\mathbb{F}}
\newcommand{\A}{\mathbb{A}}
\newcommand{\PP}{\mathbb{P}}
\newcommand{\E}{\mathbb{E}}
\newcommand{\eps}{\varepsilon}
\newcommand{\dd}{\,\mathrm{d}}
\begin{document}
\section*{UnsolvedMath ID: 20000830}
\textbf{Source problem code:} catalogue entry
\subsection{Definitions and mathematical statement}
For a fixed Hilbert polynomial $P$, the source asks whether the moduli space $\operatorname{BrV}_P(\PP^n)$ of equidimensional branchvarieties connected in codimension one is connected whenever nonempty. A branchvariety here means a reduced proper scheme $X$ with a finite morphism $f:X\to\PP^n$, with $P(m)=\chi(X,f^*\mathcal O(m))$. The field and stack/coarse-space convention are not fixed by the extracted wording.

The attached aid classifies the special case of degree-two branchcurves over $\PP^1_k$, with $k$ algebraically closed of characteristic different from $2$ and
\[
P(m)=2m+1-g.
\]
For $g\ge0$, write $L=\mathcal O_{\PP^1}(g+1)$ and $V_g=H^0(\PP^1,L^{\otimes2})$. Trace splitting presents the finite algebra as $\mathcal O\oplus L^{-1}$ with multiplication specified by a section $b\in V_g$; geometric reducedness corresponds to $b\ne0$. The aid identifies the stack as
\[
[(V_g\setminus\{0\})/\mathbb G_m],\qquad c\cdot b=c^2b.
\]
It is smooth and irreducible and is a $\mu_2$-gerbe over its coarse space $\PP^{2g+2}$. No triviality of this gerbe is asserted.

For $g\ge0$ every such source is connected, hence for these pure curves connected in codimension one. The connectedness restriction is therefore automatic. Given two section representatives, choose an auxiliary section linearly independent of each when necessary; the two affine lines interpolating through it avoid the zero section. The aid's resulting chain of at most two $\mathbb A^1$-families connects any two geometric points while all fibers remain reduced, connected and of the same Hilbert polynomial.

At $g=-1$, the complete degree-two reduced stack is $B\mu_2$, represented geometrically by the disconnected split double cover, and the connected-in-codimension-one locus is empty. Thus this exception does not contradict the original ``when nonempty'' question. The classification and deformation bridge do not prove connectedness for arbitrary $n$, degree, dimension or Hilbert polynomial.
\subsection{Short English statement}
Ask whether a nonempty fixed-polynomial branchvariety moduli space is connected. The supplied aid answers a degree-two curve case by an explicit quotient stack and deformation chains, with a separate disconnected-source exception at genus $-1$.
\subsection{Sources}
\begin{itemize}
\item Local UnsolvedMath registry, permanent ID $20000830$, source code catalogue entry; canonical \texttt{aim-arithmetic-geometry-notes.json}, \texttt{notes[75]}.
\item American Institute of Mathematics, workshop ``Components of Hilbert schemes,'' problem 29: \url{https://aimath.org/WWN/hilbertschemes/hilbertschemes.pdf}. The notation is OCR-rendered but the question is marked exact. Original literature and remarks fields are empty; no PDF inspection is claimed here.
\item Attached research attempt 1: partial result, recorded run status open, rigor proved, novelty confidence low; registry status partially solved, classification PARTIAL-PROGRESS. Its classification is machine generated and not peer reviewed; independently verify before citation. No authoritative bibliography for the quotient description is supplied.
\end{itemize}
\end{document}
